\documentclass[book.tex]{subfiles}
\begin{document}

\chapter{Geometric Convolutional Networks}

\label{chap:geometric_convolutional}
\noindent\rule{11cm}{0.4pt} \\
\textbf{Prerequisites:} \autoref{chap:mpnn} \\
\textbf{Difficulty Level:} ** \\
\noindent\rule{11cm}{0.4pt}
\vspace{5mm}

"For many industries, simulation is now becoming a guide to design rather than a mere final validation tool. A classical approach to reducing the computational complexity is to use surrogate modelling via Gaussian Process (GP) regressors, or others, trained to interpolate the performance landscape given a low dimensional parametrization of the design space. However, the regressors sometimes offer poor performance and are specific to a particular parameterization. Like other learning-based methods, Geometric Convolutional Neural Networks can be used to build surrogate models of numerical solvers. However, it suffers none of the drawbacks of previous surrogate methods, such as the ones mentioned above. It is agnostic to the shape parameters as it processes directly the mesh representation of the design. Hence, optimization or design parameters are decoupled from the learning problem and a single predictor can be trained with a large amount of data and used for many optimization tasks. In this presentation we will present two methods of learnable parametrization for geometries that eliminate the drawbacks of standard shape parameterization techniques for optimisations. Finally, we will briefly introduce Neural Concept Shape (NCS), a software platform that lets engineering teams, at all levels of expertise, implement neural network technologies into their workflows. The NCS user is working in a fully managed environment, where everything has been thought to make his work, at the frontier between data-science and traditional engineering, as effective and efficient as possible."

References \cite{baque2018geodesic}, \cite{remelli2020meshsdf}, \cite{durasov2021masksembles}.
\end{document}